% =====================================================================
%  ХУРААНГУЙ (монгол) ба ABSTRACT (англи)
% =====================================================================
\chapter*{ХУРААНГУЙ}
\addcontentsline{toc}{chapter}{Хураангуй}

Энэхүү төслийн ажилд кофе шопын үйлчилгээний урсгалыг автоматжуулсан
ухаалаг захиалгын системийг NFC стикертэй хамт нэг цогц бүтээгдэхүүн
болгон боловсруулж, ажиллагаатай кофе шопуудад шууд борлуулах
бизнесийн төслийн баримт бичгийг бүрэн хэмжээгээр гаргав. Систем нь
зочин ширээн дээр байрлах NFC стикерийг гар утсаар уншуулснаар цахим
цэс автоматаар нээгдэж, захиалгаа өгч, төлбөрийг нэг цогц орчинд
төлснөөр захиалгаа ширээндээ хүргүүлж авах боломжийг олгоно.

Төслийн бизнесийн загвар нь үйлчилгээ үзүүлэх бус, бүтээгдэхүүнийг шууд
борлуулах зарчимд суурилна: кофе шоп бүрд ширээ тус бүрт нэг удаа NFC
стикер байршуулж, тухайн кофе шопын цэс болон холбогдох системийг
тохируулан өгснөөр ханган нийлүүлэлтийн ажил дуусна. Дараа нь кофе
шоп өөрийн админ эрхээр нэвтрэн цэсээ шинэчлэх -- шинэ бүтээгдэхүүн
нэмэх, хасах, үнийн өөрчлөлт оруулах боломжтой. Ингэж олон кофе шоптой
давтан ажиллаж, нэгж борлуулалт бүрээс ашиг олно.

Судалгааны хүрээнд кафе-ийн уламжлалт үйлчилгээний урсгалын бэрхшээлүүдийг
тодорхойлж, NFC, QR код, өөрөө үйлчлэх терминал (kiosk), аппликешнд
суурилсан шийдлүүдийг харьцуулан шинжилэв. Харьцуулалтаас үзвэл NFC стикер
нь ширээг автоматаар таних, нэг алхамаар цэс нээх, гадна орчны нөлөөнд
тэсвэртэй байх зэрэг давуу талыг өгөх бөгөөд орчин үеийн Android болон
iPhone утаснууд стикерийг арын дэвсгэртэд (background) унших чадвартай
тул тусгай аппликешн таталгүйгээр ашиглах боломжтойг харуулав.

Төслийн баримт бичигт төслийн танилцуулга, нөхцөл байдлын шинжилгээ,
маркетингийн төлөвлөгөө (нэг удаагийн багц үнийн бүтэцтэй), удирдлага,
зохион байгуулалт болон хүний нөөцийн төлөвлөгөө, системийн загвар,
санхүүгийн төлөвлөгөө, эрсдэлийн үнэлгээг багтаав. Системийн загварын
хүрээнд гурван давхаргат (клиент -- аппликешн -- өгөгдөл) архитектур,
NFC уншилтаас төлбөр баталгаажуулалт хүртэлх урсгалын диаграмм,
өгөгдлийн сангийн объект-харилцааны (ER) загвар, REST API-ийн
тодорхойлолт, аюулгүй байдлын загвар, туршилтын төлөвлөгөөг
багтаав. Хөгжүүлэлтийн технологиор React (PWA), NestJS, PostgreSQL,
Redis болон Docker-ийг сонгосон. Санхүүгийн урьдчилсан тооцоогоор
нийт хөрөнгө оруулалт 55.5 сая төгрөг буюу 3 жилийн дотор 171\%-ийн
хөрөнгө оруулалтын өгөөжтэй (ROI), ойролцоогоор 35 сарын дотор
бүрэн нөхөгдөхөөр байна.

\noindent\textbf{Түлхүүр үгс:} NFC, ухаалаг захиалгын систем, PWA,
REST API, төлбөрийн интеграци, бүтээгдэхүүний борлуулалт.

\chapter*{ABSTRACT}
\addcontentsline{toc}{chapter}{Abstract}

This project report presents the full business documentation for a
smart ordering system for coffee shops, developed together with an
NFC sticker as a single product and sold directly to operating coffee
shops. The system allows a customer to tap their phone on an NFC
sticker attached to the table, which opens a web-based menu with the
table automatically identified, place an order, and pay online so
that the order is delivered directly to their table --- without
standing in line at the counter.

The business model of the project is product-based direct sales
rather than a service: for each coffee shop, one NFC sticker is
installed on every table and the menu and related system are
configured for that shop, completing the delivery. Afterwards, the
coffee shop signs in with its own administrator account to update
its menu --- adding or removing products and changing prices. The
project earns profit by repeating this delivery across many coffee
shops.

The report analyses the shortcomings of the traditional counter-service
workflow and compares NFC-based identification with QR codes,
self-service kiosks and dedicated mobile applications. The analysis shows
that NFC stickers provide one-tap access to the menu, reliable table
identification, and high resilience to environmental conditions, while
modern Android and iPhone devices support background tag reading without
installing a dedicated application.

The document contains the project introduction, situation analysis,
marketing plan with a one-time package pricing structure, management
and human resources plan, system design, financial plan, and risk
assessment. The system design covers a three-tier architecture
(client --- application --- data), a diagram of the NFC-based order
and payment flow, an entity-relationship model of the database, a REST
API specification, a security design, and a test plan. React (PWA),
NestJS, PostgreSQL, Redis and Docker were selected as the
implementation technology stack. According to the preliminary financial
estimates, the total investment is 55.5 million MNT, with a three-year
return on investment (ROI) of 171\% and a full payback period of
approximately 35 months.

\noindent\textbf{Keywords:} NFC, smart ordering system, PWA, REST API,
payment integration, product sales.
